\section{A brief overview of the regional climate}

In \autoref{fig:Temperature_2012_2021_mean}--\ref{fig:SeaIce_2012_2021_mean} one could observe how the regional climate changes throughout the year. By observing the pressure one can see that extratropical cyclones are more common during the winter months, due to the large pressure gradient. One can observe that cyclones are most prominent near the Icelandic and Aleutian lows. However, according to \citet{serreze2014} although more cyclones occur in the storm tracks cyclones in the Arctic ocean are by no means a rare phenomena. The winter is also associated by large cryosphere extent and very cold temperatures. This can easily be explained by the polar night, leading to the only heat source being the ocean, which mostly is covered by sea ice, leading to a very cold and dry climate. Somewhat warmer temperatures can be observed during the spring, although the cryosphere extent is still large. During the summer months, the temperatures are above zero, indicating melting of the ice cap. The sea ice extent is thus at a minimum during the autumn equinox \citep{serreze2014}.  

